    \begin{question}{0.1}
        (20分) (1) 证明：对于任意二阶张量 \( \sigma_{ij} \)，其偏张量 \( s_{ij} = \sigma_{ij} - \frac{1}{3}\sigma_{kk}\delta_{ij} \) 的迹为零，即 \( s_{kk} = 0 \)。

        (2) 已知应力张量的三个主不变量为：
        \[
        I_1 = \sigma_{kk}, \quad I_2 = \frac{1}{2}(\sigma_{ii}\sigma_{jj} - \sigma_{ij}\sigma_{ji}), \quad I_3 = \det(\sigma_{ij})
        \]
        证明应力偏量$s_{ij}$的第二、第三不变量 \( J_2, J_3 \) 与应力不变量的关系为：
        \[
        J_2 = \frac{1}{3}I_1^2 - I_2
        \]

    \end{question}
    
    \begin{solution}
        (1) 由偏张量定义：
        \[
        s_{ij} = \sigma_{ij} - \frac{1}{3}\sigma_{kk}\delta_{ij}
        \]
        求迹（令 \( i = j \) 并求和）：
        \[
        s_{kk} = \sigma_{kk} - \frac{1}{3}\sigma_{kk}\delta_{kk}
        \]
        注意 \( \delta_{kk} = \delta_{11} + \delta_{22} + \delta_{33} = 3 \)，代入得：
        \[
        s_{kk} = \sigma_{kk} - \frac{1}{3}\sigma_{kk} \cdot 3 = \sigma_{kk} - \sigma_{kk} = 0
        \]
        证毕。这表明偏张量仅描述畸变部分，不含体积变化。

        (2) 证明：
定义应力偏量 \( s_{ij} = \sigma_{ij} - \frac{1}{3}I_1\delta_{ij} \)。

首先证明 \( J_2 \) 的关系：
由 \( J_2 = \frac{1}{2}s_{ij}s_{ij} \)，代入 \( s_{ij} \) 表达式：
\[
\begin{aligned}
2J_2 &= s_{ij}s_{ij} = \left(\sigma_{ij} - \frac{1}{3}I_1\delta_{ij}\right)\left(\sigma_{ij} - \frac{1}{3}I_1\delta_{ij}\right) \\
&= \sigma_{ij}\sigma_{ij} - \frac{2}{3}I_1\sigma_{ij}\delta_{ij} + \frac{1}{9}I_1^2\delta_{ij}\delta_{ij} \\
&= \sigma_{ij}\sigma_{ij} - \frac{2}{3}I_1^2 + \frac{1}{3}I_1^2 \quad (\text{因为 } \sigma_{ij}\delta_{ij}=I_1, \delta_{ij}\delta_{ij}=3) \\
&= \sigma_{ij}\sigma_{ij} - \frac{1}{3}I_1^2
\end{aligned}
\]
注意 \( I_2 = \frac{1}{2}(I_1^2 - \sigma_{ij}\sigma_{ij}) \)，所以 \( \sigma_{ij}\sigma_{ij} = I_1^2 - 2I_2 \)。代入：
\[
2J_2 = (I_1^2 - 2I_2) - \frac{1}{3}I_1^2 = \frac{2}{3}I_1^2 - 2I_2
\]
    \end{solution}
